\documentclass[11pt]{article}
\linespread{1.5}
\usepackage{fullpage}
\usepackage{amsmath,theorem,amssymb,graphicx,pgfplots,tabularx,placeins}
\usepackage[semicolon,authoryear]{natbib}
\usepackage{caption}
\usepackage{subcaption}
\usepackage{csquotes}
\usepackage{epstopdf}
\pgfplotsset{compat=1.18}

\newcommand{\lb}{\label}
\newtheorem{thm}{Theorem}
\newtheorem{prop}{Proposition}
\newtheorem{definition}{Definition}

\newcommand{\bit}{\begin{itemize}}
	\newcommand{\eit}{\end{itemize}}
\newcommand{\ben}{\begin{enumerate}}
	\newcommand{\een}{\end{enumerate}}
\newcommand\setItemnumber[1]{\setcounter{enumi}{\numexpr#1-1\relax}}

\title{Retake Exam Ph.D. Macroeconomics II \\
Department of Economics, Uppsala University\\
August 20, 2026}
\author{}
\date{}

\begin{document}
\maketitle

\section*{Instructions}
\bit
	\item Writing time: 5 hours.
	
	\item The exam is closed book. 
	
	\item The exam has 70 points in total 
	
	\item A passing grade requires a) at least 30 points on  the exam, and b) 50 points in total for the course (incl the points you have from your problem sets; the problem set points will be rescaled so that they sum to 30 points: part I 25\%, part II 75\% of problem set points).
	
	\item Start each question on a new paper. Write your anonymous code on all answer pages.
	
	\item You may write your solutions by pen or pencil; use your best handwriting.
	
	\item Answers shall be given in English.

	\item Motivate your answers carefully; if you think you need to make additional assumptions to answer the questions, state them.
	
	\item If you have any questions during the exam, you may call Erik \"{O}berg (+46 730 606 796) or Teodora Borota (+46 739 262 330) at any time between 3 PM and 5 PM.
\eit

\newpage

\section{Question 1: Part I (20 points)}

\textit{[Part I question to be inserted (20 points).]}

\newpage

\section{Investment adjustment costs in the RBC model (20 points)}

Consider the RBC model studied in class in which firms own the capital stock and face convex investment adjustment costs. The cost function is
\begin{eqnarray}
C(I_t,K_t)=\frac{\phi}{2}\left(\frac{I_t}{K_t}-\delta\right)^2 K_t, \qquad \phi\geq 0. \nonumber
\end{eqnarray}
A representative firm chooses labor, investment, and next period's capital to maximize
\begin{eqnarray}
\max_{N_t, I_t, K_{t+1}} && \mathbb{E}_0 \sum_{t=0}^{\infty} M_{0,t} \left(A_t F(K_t, N_t) - W_t N_t - I_t - C(I_t,K_t)\right) \nonumber \\
\text{s.t.} && K_{t+1} = I_t + (1-\delta)K_t, \nonumber
\end{eqnarray}
where $M_{0,t}$ is the stochastic discount factor, $A_t$ is exogenous TFP, $F$ has constant returns to scale, and $\delta\in(0,1)$ is the depreciation rate. Let $q_t$ denote the Lagrange multiplier on the capital accumulation constraint (marginal $q$).

\ben
\item Write the Lagrangian for the firm's problem and derive the first-order conditions for $N_t$, $I_t$, and $K_{t+1}$. 

\item Using the quadratic cost function above, show that optimal investment satisfies
\begin{eqnarray}
\frac{I_t}{K_t}=\frac{1}{\phi}(q_t-1)+\delta. \nonumber
\end{eqnarray}

\item Consider a deterministic steady state. Using the first-order conditions, show that $\bar q=1$, $\bar I=\delta\bar K$, and $C(\bar I,\bar K)=0$. What happens to the steady-state level of investment as $\phi$ becomes larger?

\item Consider a transitory AR(1) positive TFP shock. Qualitatively describe how a \emph{larger} $\phi$ changes the dynamic responses of $q_t$, investment, and output.
\een

\newpage

\section{Hiring subsidies to reduce unemployment (15 points)}

Consider the standard continuous-time DMP model of frictional unemployment with Nash bargaining and exogenous separations, as studied in class. A legislator is considering giving a constant subsidy $s$ to any firm that employs a worker in order to promote job creation and reduce the economy's unemployment rate. With the subsidy, the profit flow of a producing firm is $y+s-w$, where $y$ is match output and $w$ is the wage. The legislator is considering two options to finance this subsidy under a balanced budget: (i) by levying a labor income tax $\tau_e$ on currently employed households, or (ii) by levying an income tax $\tau_u$ on currently unemployed households (the second financing option may also be interpreted as a reduction in unemployment benefits).

\ben
\item Will the hiring subsidy be successful in reducing the unemployment rate? Does it depend on how it is financed?

\item Suppose the policy is successful in reducing the economy's unemployment rate. Is it clear that the policy has also raised social welfare?
\een

\newpage

\section{Employment risk in the Aiyagari model (15 points)}

Consider a stationary Aiyagari economy with a continuum of ex-ante identical households. Each household's employment status $\epsilon_t\in\{e,u\}$ follows a two-state Markov process: when employed ($\epsilon=e$), the household supplies one unit of labor; when unemployed ($\epsilon=u$), the household supplies zero labor. When employed, labor income is the wage $w$; when unemployed, the household receives benefit $b$. Let $e$ and $u=1-e$ denote the invariant masses of employed and unemployed households. Households can save in a risk-free asset $a'$ (which in equilibrium sums up the capital stock $K$) but cannot borrow. Their budget constraint is
\begin{eqnarray}
	c+a'=w\cdot\mathbf{1}_{\{\epsilon=e\}}+b\cdot\mathbf{1}_{\{\epsilon=u\}}+(1+r)a, \nonumber
\end{eqnarray}
Preferences are
\begin{eqnarray}
E_0\sum_{t=0}^{\infty}\beta^t u(c_t), \nonumber
\end{eqnarray}
where $u$ satisfies standard assumptions, including prudence. Firms are competitive and produce with
\begin{eqnarray}
Y=K^\alpha L^{1-\alpha}, \nonumber
\end{eqnarray}
with depreciation rate $\delta$. Aggregate labor used in production equals the mass of employed households, $L=e$.

For parts 1--3 below, treat the unemployment benefit $b\geq 0$ as an exogenous transfer (financing is abstracted from).

\ben
\item Write the household's and firms' problems and derive their respective first-order conditions.

\item Define a stationary competitive equilibrium.

\item Draw and explain the Aiyagari diagram, with aggregate asset supply $A^s(r)$ and capital demand $K^d(r)$. Why must the equilibrium interest rate satisfy $r<1/\beta-1$?

\item Suppose now that unemployment benefits are financed by a proportional tax $\tau$ on employment income. The household budget constraint is then
\begin{eqnarray}
c+a'=(1-\tau)w\cdot\mathbf{1}_{\{\epsilon=e\}}+b\cdot\mathbf{1}_{\{\epsilon=u\}}+(1+r)a, \nonumber
\end{eqnarray}
and the government's balanced-budget condition in steady state is
\begin{eqnarray}
\tau w\, e = b\, u. \nonumber
\end{eqnarray}
Show that if the government raises benefits $b$, steady-state output is lower. Describe the mechanism in words.
\een

\end{document}
